\section{Auswertung Netz Schwanderbärgli optimiert V1.1}
\label{sec:auswertung netz schwanderbärgli optimiert v1.1}
Im Anschluss an die Datenaufbereitung wurden die Messdaten der
Feldkampagne 2026 mit dem Programm GeoSuite (swisstopo LTOP)
ausgewertet. Die bereinigten Beobachtungen wurden im `.mes`-Dateiformat
und die Näherungskoordinaten der Netzpunkte im `.koo`-Dateiformat in
LTOP eingelesen. Als Grundlage für die Näherungskoordinaten dienten
die Koordinaten aus dem Jahr 2022. Für die im Rahmen der Feldkampagne
2026 mittels GNSS bestimmten Punkte wurden die aktuellen
GNSS-Koordinaten übernommen. Die Netzausgleichung erfolgte nach der
Methode der kleinsten Quadrate (MdKQ). Dabei wird angenommen, dass die
Beobachtungen stochastisch unabhängig und die Messabweichungen
normalverteilt sind.

Im Rahmen dieses Arbeitspakets wird das Netz Schwanderbärgli in der
Variante V1.1 «optimiert» ausgewertet. Ziel ist es, durch eine geeignete
Auswahl und Kombination der vorhandenen tachymetrischen Messungen (TPS),
GNSS-Beobachtungen, Nivellementdaten und trigonometrischen
Höhenbestimmungen eine möglichst einfache Netzkonfiguration zu
entwickeln, ohne die erforderliche Genauigkeit und Zuverlässigkeit
wesentlich zu beeinträchtigen. Dazu werden verschiedene
Netzkonfigurationen und Lagerungsarten untersucht und hinsichtlich
ihrer Koordinatengenauigkeiten, Beobachtungsverbesserungen und
Zuverlässigkeitskennwerte miteinander verglichen. Anhand dieser
Ergebnisse soll beurteilt werden, welche Beobachtungen für eine
zukünftige, vereinfachte Überwachung des Schwanderbärgli weiterhin
erforderlich sind.

\subsection{Provisorischer Abriss TPS, GNSS und Nivellement}
\label{sec:provisorischer abriss tps gnss und nivellement}
Als erster Schritt wurde ein provisorischer Abriss erstellt, um eine erste
Qualitätskontrolle des Netzes durchzuführen. Beim provisorischen Abriss werden
alle Punkte als Festpunkte behandelt. Es handelt sich somit um eine reine
Plausibilitätskontrolle der gemessenen Richtungen und Distanzen anhand der
bekannten Näherungskoordinaten.

Der provisorische Abriss zeigte insgesamt ein plausibles und stimmiges Messnetz.
Die Hin- und Rückmessungen der Distanzen stimmen sehr gut überein (grösste
Differenz 5 mm bei 41 Distanzpaaren), und auch die gegenseitig gemessenen
Höhenwinkel zeigen im Median keine Abweichungen. Es wurden keine groben Fehler
in den Beobachtungen festgestellt. Einzelne Beobachtungen zeigten grössere
Abweichungen. Diese treten gruppiert auf, etwa als gemeinsame Verschiebung der
Distanzen ab Station 2000, und sind durch Bewegungen im Messgebiet sowie durch
die verschiedenen Näherungskoordinaten (GNSS Koordinaten 2026 und Koordinaten 2022) erklärbar.
Sie blieben somit im erwarteten Bereich eines provisorischen Abrisses. Die
Auswertung ergab 360 Lagebeobachtungen (davon 166 Richtungen, 166 Distanzen und
28 GNSS-Koordinaten) sowie 183 Höhenbeobachtungen (davon 169 Höhendifferenzen
und 14 GNSS-Höhen) bei 118 Stationen.

Eine Auffälligkeit zeigte die gegenseitige Höhenkontrolle zwischen den Punkten
510 und EX9500. Dort beträgt die Differenz der Hin- und Rückmessung 53.5 mm bei
einer Distanz von rund 1264 m. Da zum aktuellen Zeitpunkt noch trigonometrische
Höhenbestimmungen fehlen, wird diese Beobachtung vorerst nicht weiter beurteilt.
Das Problem wird zu einem späteren Zeitpunkt angeschaut.

\subsection{Weiche Lagerung TPS, GNSS und Nivellement}
\label{sec:weiche lagerung tps gnss und nivellement}
Nach dem provisorischen Abriss wurde das Netz in einer weichen Lagerung
ausgeglichen. Auf eine vorgängige freie Netzausgleichung wurde verzichtet.
Stattdessen wurden die vorhandenen Referenzkoordinaten bereits in der
ersten Ausgleichung berücksichtigt. Durch die weiche Lagerung können
diese innerhalb ihrer vorgegebenen Genauigkeiten angepasst werden,
ohne dass das Netz durch eine starre Festhaltung der Referenzpunkte
unnötig gezwängt wird. Dies ist insbesondere bei einem ausgedehnten
Überwachungsnetz mit unterschiedlichen Beobachtungstypen und
Referenzkoordinaten aus verschiedenen Messepochen von Bedeutung.

Ziel der weichen Lagerung war es, die innere Konsistenz der
Beobachtungen zu überprüfen, auffällige Messungen zu identifizieren
und eine geeignete Ausgangslage für die weitere Netzoptimierung
zu schaffen. Dazu wurden insbesondere die Verbesserungen der
Beobachtungen sowie die Genauigkeits- und Zuverlässigkeitskennwerte
der Ausgleichung untersucht.
Die Optimierung erfolgte schrittweise in den Varianten
weiche Lagerung v0 bis v5. Dabei wurden auffällige Beobachtungen nicht
allein aufgrund der statistischen Kennwerte eliminiert, sondern
zusätzlich anhand der Messgeometrie und der bekannten Bedingungen
während der Feldmessung beurteilt.

\subsubsection{V0}
\label{sec:v0}
